\section{The character extension and an explicit imaginary quadratic test}
\label{sec:characters}

\subsection{Conventions and the part already supplied by F-SM-13}

This section treats task~1(d).  The primitive Dirichlet numerical
formula is ALPOGE\_BRIDGES F-SM-13, whose proof is not repeated.
The new work is the precise imaginary quadratic specialization,
an ideal-indexed prime-carrying map with its norm obstruction, and
an admissible finite-$S$ example.  None of these arguments establishes
the missing zero-location-free narrow-strip transfer in Masullo
\S12.3, p.~514, or the invocation of that transfer in the proof of
Theorem~12.5.1.1, p.~520.

For $h\in C_c^\infty(\mathbb R_{>0})$ with $h(1)=0$, write
\[
 \widehat h(z)=\int_0^\infty h(t)t^z\,\frac{dt}{t},\qquad
 H(s)=\widehat h(s-\tfrac12),\qquad
 P(h)=\widehat h(\tfrac12)+\widehat h(-\tfrac12).
 \tag{MA-CH.1}
\]
For a primitive Dirichlet character $\chi$ of conductor $q$ and
parity $\chi(-1)=(-1)^a$, F-SM-13 gives
\[
 \begin{split}
 A_a(h)&=\int_0^\infty
   \frac{e^{-(a+1/2)x}}{1-e^{-2x}}
       \bigl(h(e^x)+h(e^{-x})\bigr)\,dx,\\
 B_\chi(h)&=\sum_{n\ge1}\frac{\Lambda(n)}{\sqrt n}
       \bigl(\chi(n)h(n)+\overline\chi(n)h(n^{-1})\bigr),\\
 A_a(h)+B_\chi(h)&=\boldsymbol1_{q=1}P(h)
    -\sum_{\rho}m_\chi(\rho)\widehat h(\rho-\tfrac12).
 \end{split}
 \tag{MA-CH.2}
\]
Here and below the full sum is over the nontrivial completed zeros;
trivial finite-$L$ zeros are not counted separately.  The last
identity with the full sum is unconditional.  The use of $A$ in
F-SM-13 is exactly the replacement of that sum by the sum in
$|\operatorname{Re}\rho-1/2|<\alpha$, for every $\alpha>0$.
The conductor and parity are primitive ones: an imprimitive
Dirichlet series with deleted Euler factors is a different
function.  The conjugate character on $h(n^{-1})$ is essential.
For $q>1$ there is no polar term.  These are the same conventions
as F-SM-13, equations preceding and including
\texttt{SM-numeric-strip}, and its Mellin-inversion calculation.

To fix the adelic orientation for both fields, use GRH\_NIGHT
MT-M22 and MT-M24.  If the primitive finite Euler factors are
$(1-c_{\mathfrak p}(N\mathfrak p)^{-s})^{-1}$, let $\beta$ be
Meyer's idele-class character specified by
\[
       \beta(p_v)^{-1}=c_{\mathfrak p_v}
       \quad(v\notin S),\qquad
       \omega_s=\beta|\cdot|^s .
 \tag{MA-CH.3}
\]
The inverse in this equation is Meyer's inverse in the Euler
multiplier, not a naming convention to be suppressed.  Put
$C=C_K$, $C^1=\ker|\cdot|$, and $\xi=\beta|_{C^1}$.
Choose a norm splitting $a:\mathbb R\to C$, $|a(x)|=e^x$.
Since $\beta$ has finite order, $\beta(a(x))=1$.
With probability Haar measure $dc$ on $C^1$, the actual projector
and sector test are
\[
 \begin{aligned}
 (e_\xi f)(g)&=\int_{C^1}\xi(c)f(cg)\,dc,
       &e_\xi^2&=e_\xi,\\
 \phi_{\xi,h}(g)&=\beta(g)^{-1}h(|g|),
       &\widehat\phi_{\xi,h}(\omega_s)&=\widehat h(s).
 \end{aligned}
 \tag{MA-CH.4}
\]
The range of $e_\xi$ consists of the functions satisfying
$f(cg)=\xi(c)^{-1}f(g)$, equivalently those on which left
translation by $c$ acts by $\xi(c)$.  Substitution $u=cd$ in the
integral proves this assertion and idempotence; character
orthogonality proves the Mellin formula and the vanishing of all
other compact components.  Thus (MA-CH.4) is an actual sector
projection, not merely a spectral correspondence.  It commutes
with the adelic summation and Fourier constructions in the sense
specified by MT-M22/M24; no assertion that Fourier transform
preserves one character rather than exchanges dual sectors is
needed.

Meyer's normalization is \S3.1, author p.~16, extended to $C_K$
in \S5.1, author p.~38.  His global trace and divisor statements
are Theorems~5.8 and~5.11, author pp.~43--44 and~47.
For the actual image criterion use author pp.~46--47, and for
its transfer to the specified Masullo denominator use MT-M22
(M22.7)--(M22.10) and MT-M24 (M24.6)--(M24.9).
These are published Meyer input and the already supplied
comparison results, not consequences of $A$ or $M$.

There is a numerical normalization to retain when returning to
the draft.  Masullo Definition~A.3.2.1 and (A.3.2.1), pp.~559--560,
give compact volume
\[
 v_K=\frac{2^{r_1}(2\pi)^{r_2}h_KR_K}
                  {w_K\sqrt{|D_K|}}.
 \tag{MA-CH.5}
\]
The measure in (MA-CH.4) is his measure divided by $v_K$.
Convolution using his measure and test $\phi/v_K$ is exactly
convolution using the normalized measure and test $\phi$.
With a fixed unrescaled test the former operator, and its trace
when defined, is $v_K$ times the latter.  MT-M24 explains why
this nonzero scalar does not change the theta-image subspace;
it must nevertheless be retained in a numerical trace
comparison.  The draft's local Theorem~12.4.1.1, p.~514,
likewise divides by the norm-one local volume.  All formulas
in this section use the normalized convention.

\subsection{The imaginary quadratic explicit formula and its exact gate}

\paragraph{MA-CH-001 (the finite-order imaginary quadratic specialization).}
Let $K$ be imaginary quadratic, and let $\eta$ be a nontrivial
primitive finite-order Hecke character of conductor $\mathfrak f$.
Write $c_{\mathfrak p}=\eta(\mathfrak p)$ when $\eta$ is
unramified at $\mathfrak p$ and $c_{\mathfrak p}=0$ otherwise.
This definition uses inertia over $K$; ramification of $K/\mathbb Q$
alone does not force $c_{\mathfrak p}=0$.
Put
\[
 \begin{aligned}
 Q_\eta&=|D_K|N\mathfrak f,\\
 \Lambda_\eta(s)&=Q_\eta^{s/2}\Gamma_{\mathbb C}(s)L_K(s,\eta),\\
 \Gamma_{\mathbb C}(s)&=2(2\pi)^{-s}\Gamma(s),\\
 A_{\mathbb C}(h)&=\int_0^\infty
       \frac{e^{-x/2}}{1-e^{-x}}
          \bigl(h(e^x)+h(e^{-x})\bigr)\,dx,\\
 B_\eta(h)&=\sum_{\mathfrak p}\sum_{k\ge1}
       \frac{\log N\mathfrak p}{(N\mathfrak p)^{k/2}}
       \left(c_{\mathfrak p}^{\,k}h((N\mathfrak p)^k)
          +\overline{c_{\mathfrak p}}^{\,k}
                       h((N\mathfrak p)^{-k})\right),\\
 \tau_\eta^{\rm full}(h)&=\sum_\rho
           m_\eta(\rho)\widehat h(\rho-\tfrac12).
 \end{aligned}
 \tag{MA-CH.6}
\]
Then
\[
              A_{\mathbb C}(h)+B_\eta(h)
                         =-\tau_\eta^{\rm full}(h).
 \tag{MA-CH.7}
\]
The prime sum is finite, the integral and zero sum converge
absolutely, and this is an unconditional equality of continuous
linear functionals on the indicated test space.  For the trivial
character the same statement has $P(h)$ on the right before
$-\tau^{\rm full}$.

For the open-strip distribution
\[
 \tau_{\eta,\alpha}(h)=
     \sum_{|\operatorname{Re}\rho-1/2|<\alpha}
            m_\eta(\rho)\widehat h(\rho-\tfrac12),
 \tag{MA-CH.8}
\]
the exact discrepancy is
\[
 A_{\mathbb C}(h)+B_\eta(h)+\tau_{\eta,\alpha}(h)
       =-\!\sum_{|\operatorname{Re}\rho-1/2|\ge\alpha}
            m_\eta(\rho)\widehat h(\rho-\tfrac12).
 \tag{MA-CH.9}
\]
Consequently (MA-CH.7) with every positive narrow strip follows
from RH for this one $L_K(s,\eta)$, hence from $M$.  The
corresponding Dirichlet statement follows from $A$.  The
imaginary quadratic assertion does not follow from the
Dirichlet scope of $A$.

\emph{Proof of the new specialization.}
A finite-order continuous character is trivial on the connected
complex archimedean group; thus its archimedean factor is the
single $\Gamma_{\mathbb C}(s)$ in (MA-CH.6).  The primitive
completion is entire for nontrivial $\eta$, and its functional
equation is
$\Lambda_\eta(s)=W_\eta\Lambda_{\bar\eta}(1-s)$.
The full explicit formula, equivalently Meyer's
Theorems~5.8/5.11 on the sector (MA-CH.4), is unconditional.
To evaluate its terms with the stated normalization, the same
contour pairing as in F-SM-13 is
\[
 \frac1{2\pi i}\int_{(c)}
 \left(\frac{\Lambda_\eta'}{\Lambda_\eta}(s)H(s)
       +\frac{\Lambda_{\bar\eta}'}{\Lambda_{\bar\eta}}(s)
                                                  H(1-s)\right)ds,
 \qquad c>1.
 \tag{MA-CH.10}
\]
It equals the full zero sum, with the two polar residues
subtracted only for the trivial character.  The standard
fixed-character polynomial bounds on vertical strips and the
order-one zero count, already used in MT-M18/M22 and F-SM-13,
justify the contour pairing and absolute convergence against
the rapidly decreasing Mellin transform.  Alternatively its
convergence is the summability in Meyer Theorem~5.8; no
critical-line assertion enters either justification.

The two constant logarithmic derivatives contribute
$(\log Q_\eta-2\log(2\pi))h(1)=0$.  On $\operatorname{Re}s=c$
the Euler logarithmic derivative is absolutely convergent,
and Mellin inversion gives, for $n>1$,
\[
 \frac1{2\pi i}\int_{(c)}H(s)n^{-s}ds=n^{-1/2}h(n),\qquad
 \frac1{2\pi i}\int_{(c)}H(1-s)n^{-s}ds=n^{-1/2}h(n^{-1}).
 \tag{MA-CH.11}
\]
Hence the finite terms in (MA-CH.10) are exactly $-B_\eta(h)$,
including the conjugate coefficient on the inverse argument.
For the gamma term use
\[
 \frac{\Gamma'}{\Gamma}(s)
    =-\gamma+\int_0^\infty
                    \frac{e^{-x}-e^{-sx}}{1-e^{-x}}\,dx
       \quad(\operatorname{Re}s>0).
\]
Truncate the integral, pair it with $H(s)+H(1-s)$, and apply
(MA-CH.11).  The constants pair with $h(1)=0$, while the
remaining term is exactly $-A_{\mathbb C}(h)$.  Near zero the
numerator $h(e^x)+h(e^{-x})$ is $O(x^2)$, and for large $x$
it vanishes, so the resulting integral is absolutely
convergent.  The truncated digamma expressions have a
uniform bound by a polynomial in $|\operatorname{Im}s|$ on
the fixed integration line, which permits passage to the
limit against $H$.  This proves (MA-CH.7).

The same archimedean coefficient can be checked directly from
the complex local Weil integral.  With $z=re^{i\theta}$ and
$t=|z|_{\mathbb C}=r^2$, the canonical multiplicative measure is
$dt/t\,d\theta/(2\pi)$.  The angular average
\[
 \int_0^{2\pi}\frac{d\theta}{2\pi|1-re^{i\theta}|^2}
                  =\frac1{|1-r^2|}
 \quad(r\ne1)
\]
turns that integral into
$\int_0^\infty h(t)\sqrt t\,|1-t|^{-1}dt/t$, which is
$A_{\mathbb C}(h)$ after splitting at $1$.
Thus no extra factor of $2$ is implicit in the notation
$\Gamma_{\mathbb C}$.

Equation (MA-CH.9) is subtraction of absolutely convergent
sums.  It also specifies the exact residual gate.  For a
fixed $\alpha$, its right side vanishes for every allowed
test if and only if there is no completed zero outside the
open strip.  Indeed both full and partial zero distributions
are continuous on the exponential Schwartz space and
annihilate the actual centered theta ideal.  MT-M30 allows
the norm-one vanishing restriction without loss of the
quotient, and density gives all such smooth tests.  The
zero-jet interpolation of MT-M18 then detects any one
excluded zero; the other jets can be prescribed zero.
This argument keeps multiplicities and also detects a zero
on the strip boundary.  Accordingly, equality for all
$\alpha>0$ is precisely RH for this single character, not
an additional proof of it.  For actual Hilbert operator
traces use MT-M26/MT-M27 on wide strips, or their
zero-location-dependent narrow-strip extensions.  The
displayed partial sum alone is not a claim of trace class
for a boundary-sensitive arbitrary quotient.\hfill$\square$

\subsection{An ideal-indexed transport and its escaping correction}

\paragraph{MA-CH-002 (a twisted summation embedding and a norm obstruction).}
Let $K=\mathbb Q$ with a nonprincipal primitive Dirichlet
character, or let $K$ be imaginary quadratic with the
nontrivial character $\eta$ above.  Extend its Euler
coefficient multiplicatively to integral ideals by
$c(\mathfrak a)=0$ when $\mathfrak a$ meets the conductor.
For $K=\mathbb Q$ the ideals are positive integers and
$c(n)=\chi(n)$.  If $F\in C_c^\infty(\mathbb R_{>0})$, define
\[
 (\mathcal Z_\eta F)(t)=
       \sum_{\mathfrak a\ne0}
            \frac{c(\mathfrak a)}{\sqrt{N\mathfrak a}}
                           F((N\mathfrak a)t).
 \tag{MA-CH.12}
\]
This is an injective, norm-translation-equivariant map into
the actual centered closed theta ideal identified in
MT-M22/MT-M24.  It is continuous from each fixed compact
support test space to the exponential Schwartz space.
Its Mellin transform and inverse on its image are
\[
 \begin{split}
 \widehat{\mathcal Z_\eta F}(z)
        &=L_K(z+\tfrac12,\eta)\widehat F(z),\\
 F(t)&=\sum_{\mathfrak a\ne0}
    \frac{\mu_K(\mathfrak a)c(\mathfrak a)}{\sqrt{N\mathfrak a}}
             (\mathcal Z_\eta F)((N\mathfrak a)t).
 \end{split}
 \tag{MA-CH.13}
\]
The inverse is an ideal-indexed twisted M\"obius sum.
No surjectivity onto the whole theta ideal is asserted.

Put $(U_y f)(t)=f(e^{-y}t)$ and, for a finite set $S$ of
prime ideals, put
\[
 P_{\eta,S}=\prod_{\mathfrak p\in S}
  \left(1-c_{\mathfrak p}(N\mathfrak p)^{-1/2}
                                      U_{-\log N\mathfrak p}\right).
 \tag{MA-CH.14}
\]
For $g=\mathcal Z_\eta F$ one has the exact finite-Euler formula
\[
       P_{\eta,S}g(t)=
       \sum_{\substack{\mathfrak a\ne0\\(\mathfrak a,S)=1}}
            \frac{c(\mathfrak a)}{\sqrt{N\mathfrak a}}
                                     F((N\mathfrak a)t).
 \tag{MA-CH.15}
\]
If $\operatorname{supp}F\subset[a,b]$, the correction
$P_{\eta,S}g-F$ is supported in $(0,b/q_S]$, where $q_S$
is the least norm of an omitted prime with nonzero Euler
coefficient.  Hence it eventually vanishes on each compact
subset of $\mathbb R_{>0}$ as $S$ increases through all
unramified primes.

Nevertheless fix any nontrivial completed zero $\rho$ with
$\lambda=\rho-1/2$ and $|\operatorname{Re}\lambda|<\alpha$.
For $s\ge0$ use the same centered norm as in the main section,
\[
 \|u\|_{\alpha,s}^2=
  \int_{\mathbb R}|u(e^x)|^2e^{2\alpha|x|}(1+x^2)^{s/2}\,dx.
 \tag{MA-CH.16}
\]
There is a compactly supported $F$ with $\widehat F(\lambda)=1$,
and for that $F$ every finite $S$ satisfies
\[
 \|P_{\eta,S}\mathcal Z_\eta F-F\|_{\alpha,s}
   \ge
   \left(\int_{\mathbb R}
        e^{2\operatorname{Re}\lambda x-2\alpha|x|}
                          (1+x^2)^{-s/2}\,dx\right)^{-1/2}>0.
 \tag{MA-CH.17}
\]
This excludes convergence in this norm, despite eventual
agreement on every compact radial set.  The statement is
unconditional whenever the strip contains the chosen zero;
wide strips contain every nontrivial zero.  Applying it on
every arbitrarily narrow strip uses $A$ in the Dirichlet case,
or the single-character RH assumption in the Hecke case.

\emph{Proof.}
For each $t>0$, only the finitely many ideals with norm at
most $b/t$ occur in (MA-CH.12), locally uniformly in $t$.
Termwise differentiation is therefore valid.  Mellin
integration is absolutely convergent for
$\operatorname{Re}z>1/2$ and gives the first equation in
(MA-CH.13).  The finite Hecke $L$-function is entire in
the nontrivial cases considered here.  On each fixed
vertical strip it has at most polynomial growth; multiplication
by the compact-test Mellin transform gives rapid decay of
every desired order.  Mellin inversion may consequently
be shifted to any line $\operatorname{Re}z=-N$.
It gives $D^j\mathcal Z_\eta F(t)=O_{N,j}(t^N)$ as
$t\downarrow0$, with $D=t\,d/dt$, while the function is
identically zero for $t>b$.  The same argument bounds every
exponential Schwartz seminorm by finitely many seminorms of
$F$ on its fixed support.  Thus $g(e^x)$ belongs to the
exponential Schwartz space and depends continuously on $F$.

At every completed nontrivial zero $\rho$, the first equation
in (MA-CH.13) vanishes to at least its full multiplicity at
$z=\rho-1/2$.  MT-M22 (M22.10), followed by the centering
in MT-M24 (M24.3)/(M24.9), identifies exactly these simultaneous
jet vanishings with the actual centered theta ideal.
This proves membership using the specified adelic quotient,
without replacing its denominator by a merely formal
Euler-product ideal.

The right side of the second equation in (MA-CH.13) is again
finite for each fixed $t$: $g(u)=0$ for $u>b$.
Substituting (MA-CH.12) gives a finite double sum.  At a
nonunit ideal $\mathfrak b$ its coefficient is
\[
   \frac{c(\mathfrak b)}{\sqrt{N\mathfrak b}}
                         \sum_{\mathfrak a\mid\mathfrak b}
                                           \mu_K(\mathfrak a)=0;
\]
at the unit ideal it is $1$.  This proves the inverse,
including ideals meeting the conductor, where the extended
coefficient is zero.  In particular the map is injective.
The radial scaling in every summand proves translation
equivariance directly.

The same finite cancellation with only primes in $S$ proves
(MA-CH.15).  A nonunit ideal with nonzero coefficient and
no prime factor in $S$ has norm at least $q_S$, which proves
the support assertion.  There are only finitely many primes
of bounded norm, so $q_S\to\infty$ along any exhaustion.
Every factor in (MA-CH.14) preserves the theta ideal, either
by translation invariance or by its entire Mellin multiplier.

The functional $u\mapsto\widehat u(\lambda)$ is continuous
for the norm (MA-CH.16), with dual norm equal to the square
root of the integral in (MA-CH.17), by Cauchy--Schwarz.
Strict interiority of $\lambda$ makes that integral finite.
It vanishes on the theta ideal and its Hilbert closure.
Choose a compact smooth $b_0$ on the log line with
$\int b_0(x)dx=1$, and set
$F(e^x)=e^{-\lambda x}b_0(x)$.  Then
$\widehat F(\lambda)=1$, whereas
$\widehat{P_{\eta,S}g}(\lambda)=0$ for every $S$.
Cauchy--Schwarz gives (MA-CH.17).  No assertion about a zero
on the boundary of the strip is used.\hfill$\square$

\paragraph{What changes in the geometric intertwiner.}
For $\mathbb Q$ with a nontrivial Dirichlet character the
compact factor in (MA-CH.4), the conductor-level finite
Schwartz data, parity at the real place, and the coefficients
$\chi(n)$ must all be included.  The trivial-character
integer theta formula is therefore not an unchanged formula
for this sector.  For imaginary quadratic $K$ its exact
adelic starting point is
\[
  e_\xi\Phi_K f(g)
      =\int_{C_K^1}\xi(c)
              \sum_{u\in K^\times} f(ucg)\,dc,
 \qquad f\in\mathcal S_0(\mathbb A_K).
 \tag{MA-CH.18}
\]
At a fixed finite level its archimedean sums are sums on
rank-two lattices coming from fractional ideals, with the
ideal-class sectors and the finite congruence data retained.
A general ray character requires its conductor data even
when the ordinary ideal class group is trivial.  To pass
to an $S$-idele presentation, $S$ must include archimedean
places, the character conductor, the differential/base-field
ramification data, and enough primes to make the relevant
$S$-ideal class group trivial; see Meyer \S5.1,
author pp.~37--38, and the hypotheses of Masullo
Theorem~12.4.1.1, p.~514.

Equations (MA-CH.12)--(MA-CH.15) construct a concrete
prime-carrying radial subspace of that actual ideal, with
ideal M\"obius inversion.  They do not assert that a single
integer theta sum parametrizes all the lattices in
(MA-CH.18), or that their inverse is bounded for a narrow
Hilbert quotient.  Equation (MA-CH.17) proves that the
finite-Euler inverses cannot supply such a norm convergence
on a strip containing a zero.  The missing step is therefore
unchanged in its analytic nature: a valid global-to-local
operator trace comparison must retain the omitted ideal-prime
terms and justify the required topology.  The smooth
identification MT-M24 alone supplies neither that limit nor
zero location.

\subsection{The cubic character, its primitive local factors, and a trace gap}

\paragraph{MA-CH-003 (the primitive conductor of the concrete cubic character).}
Take the character from MA-PF-003/MA-PF-005:
\[
 K=\mathbb Q(\zeta_3)=\mathbb Q(\sqrt{-3}),\qquad
 E=K(269^{1/3}),\qquad \eta=\psi_{269,3}.
 \tag{MA-CH.19}
\]
Either faithful character of $\operatorname{Gal}(E/K)=C_3$
has primitive conductor
\[
     \mathfrak f_\eta=(269),\qquad
     Q_\eta=3\cdot269^2.
 \tag{MA-CH.20}
\]
At the unique prime over $3$ the extension $E/K$ splits
completely, so the primitive Hecke factor is
$(1-3^{-s})^{-1}$.  At the prime over $269$ the factor is
$1$, since the character is nontrivial on inertia there.
At the unique prime $(2)$, whose norm is $4$, the character
value is $1$ and its factor is $(1-4^{-s})^{-1}$.

The non-descent and Artin induction statements are the
already proved MA-PF-003/MA-PF-005:
\[
 \eta^c=\eta^{-1}\ne\eta,
 \qquad
 L_K(s,\eta)=\frac{\zeta_{\mathbb Q(269^{1/3})}(s)}
                         {\zeta_{\mathbb Q}(s)}.
 \tag{MA-CH.21}
\]
In particular this is not the norm pullback of a Dirichlet
character.  Those background identifications are not being
reproved; the local conductor and retained Euler factors
in (MA-CH.20) are the new refinement.

\emph{Proof of the local refinement.}
At $X=-1$ the polynomial $f(X)=X^3-269$ satisfies
$v_3(f(-1))=v_3(-270)=3>2v_3(f'(-1))=2$.
Thus the strong Hensel criterion gives a root in $\mathbb Q_3$.
For an elementary verification, substitute $X=-1+9t$:
\[
       f(-1+9t)=27(27t^3-9t^2+t-10).
\]
The map $t\mapsto10+9t^2-27t^3$ is a contraction of
$\mathbb Z_3$ with Lipschitz constant at most $1/9$ and
has a fixed point.  Since $K_3$ contains $\zeta_3$, it
then contains all three roots.  The prime over $3$ splits
completely in $E/K$, proving the asserted factor and zero
conductor exponent there.

Since $269\equiv2\pmod3$, it is inert and unramified in
$K/\mathbb Q$.  Over its completion the same polynomial
is Eisenstein of degree $3$.  Its extension is totally
ramified and tame, because the residue characteristic is
$269$.  A faithful rank-one character has no inertia
invariants and has zero Swan conductor, hence Artin
conductor $1-0+0=1$ at that prime.  Outside $3\cdot269$
the polynomial discriminant $-27\cdot269^2$ is a unit,
so the splitting-field extension is unramified.  There is
no archimedean finite-order conductor at the complex place.
This proves (MA-CH.20) and the factor at $269$.

Finally $2$ is inert in $K$, so its residue field is
$\mathbb F_4$.  Modulo that prime $X^3-269$ is $X^3-1$,
whose three distinct roots are exactly $\mathbb F_4^\times$.
The prime is unramified and its Frobenius is the identity
on all three roots.  Thus $c_{(2)}=1$, as asserted.
\hfill$\square$

\paragraph{MA-CH-004 (an admissible finite-$S$ trace gap for this character).}
For the character (MA-CH.19), let $S_0$ consist of the complex
place and the primes above $3$ and $269$.
There is a real $h\in C_c^\infty(\mathbb R_{>0})$, vanishing
at $1$, such that
\[
  h(4)=1,\qquad A_{\mathbb C}(h)=0,\qquad
  B_{\eta,S_0}(h)=0,\qquad B_\eta(h)=\log2.
 \tag{MA-CH.22}
\]
Here $B_{\eta,S_0}$ restricts (MA-CH.6) to the finite
places in $S_0$.  Consequently the canonical finite-$S_0$
local Weil functional is $0$, whereas the full geometric
functional is $\log2$ and the full zero trace is $-\log2$.
The same zero trace is the narrow-strip numerical trace
under RH for this one $\eta$, hence under $M$.
The discrepancy is unconditional as a comparison with the
full trace; it does not assume $M$.

\emph{Proof.}
Choose a nonnegative smooth bump $v$ supported in
$(3.9,4.1)$ with $v(4)=1$, and a nonzero nonnegative smooth
bump $w$ supported in $(1.4,1.6)$.  Both archimedean
integrals are strictly positive.  Define
\[
             h=v-\frac{A_{\mathbb C}(v)}{A_{\mathbb C}(w)}w.
 \tag{MA-CH.23}
\]
Then $h(1)=0$, $h(4)=1$ and $A_{\mathbb C}(h)=0$.
Every integral ideal norm is an integer.  The only
prime-ideal power norm in the support of $v$ is $4$,
coming from the inert prime $(2)$ with exponent one;
there is no prime ideal of norm $2$, and no prime-power
norm in the support of $w$.  Both supports lie above $1$,
so the inverse-argument terms also vanish.
By MA-CH-003 this gives
$B_\eta(h)=\log4/\sqrt4=\log2$.
The primes in $S_0$ contribute nothing: $3^k$ is outside
the supports for every $k\ge1$, and the coefficient at
$269$ is zero.  This proves (MA-CH.22).

This $S_0$ meets the local setup's actual arithmetic
hypotheses.  It contains all field and character
ramification.  Moreover $\mathbb Z[\zeta_3]$ is
norm-Euclidean: rounding the two real coordinates in the
basis $1,\zeta_3$ gives an error $a+b\zeta_3$ with
$|a|,|b|\le1/2$ and norm
$a^2-ab+b^2\le3/4<1$.  Euclidean division follows, so
the ideal class group and its $S_0$ localization are
trivial.  The character is unramified outside $S_0$.
Thus this is not a comparison with an inadmissible
archimedean-only $S$.

The canonical complex local term is the angular integral
computed in MA-CH-001.  At an unramified finite place,
Meyer's local Weil distribution expands into the two
prime-power sums in (MA-CH.6); the normalization is
Definition~5.7/Theorem~5.8, author pp.~43--44, and
\S6.1, the displayed expansion immediately after
Definition~6.1.  At a ramified character the
remaining unit contribution is supported at norm $1$,
and the conductor/discriminant terms also multiply
$h(1)$, so they vanish here.  Accordingly its finite-$S_0$
value is $A_{\mathbb C}(h)+B_{\eta,S_0}(h)=0$.
Equation (MA-CH.7) supplies the full value and its sign.
Any transfer preserving this uncorrected finite-$S_0$
functional therefore cannot identify it with the full
global functional.  A transfer carrying the omitted
ideal-prime term is not excluded.\hfill$\square$

\subsection{Artin motives and the scope of the extension}

The character $\eta$ in (MA-CH.19) is a rank-one finite-monodromy
sheaf over $K$ away from its conductor.  Pushing it from $K$
to $\mathbb Q$ gives the rank-two Artin motive
$\operatorname{Ind}_{G_K}^{G_{\mathbb Q}}\eta$ with dihedral
image $S_3$.  Its $L$-function is (MA-CH.21), including the
local inertia-invariant factors checked in MA-CH-003.
It is irreducible, since its two characters on $G_K$ are
$\eta$ and $\eta^{-1}$, and it is odd, since complex
conjugation exchanges the two induced lines and has
eigenvalues $1,-1$.  Its analytic continuation is already
that of a Hecke $L$-function.  The unresolved narrow-strip
assertion for this motive is its own RH, not Dirichlet RH.

The precise geometric dictionary is SGA~1, Expos\'e~V,
\S7, for finite \'etale covers and finite fundamental-group
sets, followed by the idempotent-completed representation
category of Artin motives in Deligne (1979), Definition~6.1
and \S\S6.2--6.3.  Deligne (0.5.3) identifies restriction of
scalars with induction; the motives here have weight zero
and Hodge type $(0,0)$.  MA-CH-003 computes their local
conductor and inertia data; it does not assign geometric
cohomology to the analytic zero jets.  MA-CH-001 is the
adelic trace formula on a fixed character component of
$C_K$.  MA-CH-002 is a statement about the topology of its
summation ideal, and MA-CH-004 exhibits the Euler data
lost by an uncorrected local trace.  In the cote~1
trace/determinant programme, scans~079--080 and~111--112,
these are respectively a character specialization, a
failure of the proposed approximation in the required norm,
and an omitted trace term; the archive locators are carried
through ALPOGE\_BRIDGES BR-TR-01/BR-TR-02 and F-SM-11/F-SM-12,
not newly certified by a scan inspection in this section.

For the example, a field of definition for the faithful
rank-one character contains $\zeta_3$, while its induced
rank-two representation has a rational model, the standard
representation of $S_3$.  Thus the failure to descend as a
rank-one norm character is not a failure of motivic descent
of its induced object.  This distinction is exactly why
the rank-one Dirichlet scope of $A$ does not cover it.

\paragraph{Input and check ledger for this section.}
MA-CH-001 is unconditional for the full distribution; its
Dirichlet narrow specialization uses $A$, and its Hecke
narrow specialization uses the RH of that one character
(sufficiently $M$).  MA-CH-002 is unconditional with its
explicit interior-zero hypothesis and uses no outside
preprint.  MA-CH-003 and MA-CH-004 are unconditional;
using their full trace as every narrow trace again needs
the single-character RH assumption.  The Kummer example
is shared with the PF26 analysis, but its field, character,
conductor and trace-gap proof do not require the PF26
factoring preprint.  The event register records the actual
source checks, full new-proof reading, and independent
local-conductor check separately.

\paragraph{References specific to this section.}
Meyer, R. (2005). On a representation of the idele class
group related to primes and zeros of $L$-functions.
\emph{Duke Mathematical Journal, 127}, 519--595.
Author version arXiv:math/0311468v3, \S3.1, pp.~16--17;
\S5.1, pp.~37--38; Theorems~5.8 and~5.11, pp.~43--44
and~47; image criterion, pp.~46--47; \S6.1, the local
expansion immediately after Definition~6.1.
\url{https://arxiv.org/abs/math/0311468}.
Deligne, P. (1979). Valeurs de fonctions $L$ et p\'eriodes
d'int\'egrales. In \emph{Automorphic forms, representations
and $L$-functions} (Proceedings of Symposia in Pure
Mathematics, Vol.~33, Part~2, pp.~313--346).
American Mathematical Society. Equation~(0.5.3),
Definition~6.1 and \S\S6.2--6.3; the Ghate--Beineke
translation has the relevant passages on PDF pp.~4 and~20.
\url{https://mathweb.tifr.res.in/~eghate/Deligne.pdf}.
